% ============================================================================
% MUSTERLÖSUNG DS 2 - Gruppe A + B
% ============================================================================

\documentclass[11pt,a4paper]{article}

\usepackage[utf8]{inputenc}
\usepackage[T1]{fontenc}
\usepackage[ngerman]{babel}
\usepackage[left=2cm,right=2cm,top=1.8cm,bottom=1.5cm]{geometry}
\usepackage{helvet}
\usepackage{xcolor}
\usepackage{tabularx}
\usepackage{array}
\usepackage{enumitem}
\usepackage{tcolorbox}
\usepackage{fancyhdr}
\usepackage{lastpage}

\renewcommand{\familydefault}{\sfdefault}

\definecolor{spanisch}{HTML}{c41e3a}
\definecolor{grau}{HTML}{666666}
\definecolor{hellgrau}{HTML}{f5f5f5}
\definecolor{gruppeB}{HTML}{1565c0}
\definecolor{loesung}{HTML}{c62828}

\pagestyle{fancy}
\fancyhf{}
\fancyhead[L]{\small\textcolor{grau}{Spanisch Spätbeginner -- MUSTERLÖSUNG}}
\fancyhead[R]{\small\textcolor{grau}{ML-02}}
\fancyfoot[C]{\small\thepage/\pageref{LastPage}}
\renewcommand{\headrulewidth}{0.5pt}

\newcommand{\lsg}[1]{\textcolor{loesung}{\textbf{#1}}}
\newcommand{\es}[1]{\textit{#1}}

\newtcolorbox{gruppenbox}[1][]{
    colback=white,
    colframe=spanisch,
    fonttitle=\bfseries\color{white},
    colbacktitle=spanisch,
    title=#1,
    boxrule=1pt,
    arc=0pt,
    left=5pt, right=5pt, top=5pt, bottom=5pt
}

\newtcolorbox{gruppeBbox}[1][]{
    colback=white,
    colframe=gruppeB,
    fonttitle=\bfseries\color{white},
    colbacktitle=gruppeB,
    title=#1,
    boxrule=1pt,
    arc=0pt,
    left=5pt, right=5pt, top=5pt, bottom=5pt
}

\begin{document}

% ============================================================================
% KOPFBEREICH
% ============================================================================
\begin{center}
{\Large\textbf{\textcolor{loesung}{MUSTERLÖSUNG}}}\\[0.3em]
{\large DS 2 -- Freizeitorte / Wortbildung (S. 31--32 / S. 104)}
\end{center}

\vspace{0.5em}
\hrule
\vspace{1em}

% ============================================================================
% GRUPPE A
% ============================================================================
\begin{gruppenbox}[Gruppe A: ¿Adónde vamos? (S. 31--32)]

\textbf{Kasten 1: Verb ir}

\begin{tabularx}{\textwidth}{@{}|c|l||c|l|@{}}
\hline
yo & \lsg{voy} & nosotros/-as & \lsg{vamos} \\
\hline
tú & \lsg{vas} & vosotros/-as & \lsg{vais} \\
\hline
él/ella/usted & \lsg{va} & ellos/-as/ustedes & \lsg{van} \\
\hline
\end{tabularx}

\vspace{0.5em}

\textbf{Kasten 2: Freizeitorte}

la piscina = \lsg{Schwimmbad}, el cine = \lsg{Kino}, el teatro = \lsg{Theater}

el parque = \lsg{Park}, la discoteca = \lsg{Disco}, el polideportivo = \lsg{Sportzentrum}

el centro comercial = \lsg{Einkaufszentrum}, la biblioteca = \lsg{Bibliothek}, el estadio = \lsg{Stadion}

\vspace{0.5em}

\textbf{Aufgabe 2: al / a la / a los / a las (6 P.)}
\begin{enumerate}[label=\alph*), leftmargin=2em, nosep]
    \item Voy \lsg{a la} piscina todos los sábados.
    \item ¿Vamos \lsg{al} cine esta noche?
    \item Mis amigos van \lsg{al} parque después de clase.
    \item Mi hermana va \lsg{a la} biblioteca para estudiar.
    \item Vais \lsg{a las} discotecas los viernes.
    \item Los niños van \lsg{al} polideportivo.
\end{enumerate}

\vspace{0.5em}

\textbf{Aufgabe 5: Dialog Beispiellösung (8 P.)}

\lsg{A: ¡Hola! ¿Qué hacemos este fin de semana?}

\lsg{B: No sé. ¿Adónde vamos?}

\lsg{A: ¿Quieres ir al cine? Hay una película nueva.}

\lsg{B: Buena idea. ¿Y después vamos a tomar algo?}

\lsg{A: Sí, podemos ir al centro comercial.}

\lsg{B: ¡Perfecto! Quedamos el sábado a las cinco.}

\end{gruppenbox}

\vspace{1em}

% ============================================================================
% GRUPPE B
% ============================================================================
\begin{gruppeBbox}[Gruppe B: Formación de palabras (S. 104)]

\textbf{Aufgabe 1: Substantive auf -ción (8 P.)}

informar → \lsg{información}, organizar → \lsg{organización}

comunicar → \lsg{comunicación}, presentar → \lsg{presentación}

formar → \lsg{formación}, evaluar → \lsg{evaluación}

motivar → \lsg{motivación}, colaborar → \lsg{colaboración}

\vspace{0.5em}

\textbf{Aufgabe 2: Substantive auf -(i)dad (6 P.)}

responsable → \lsg{responsabilidad}, capaz → \lsg{capacidad}

posible → \lsg{posibilidad}, creativo → \lsg{creatividad}

flexible → \lsg{flexibilidad}, productivo → \lsg{productividad}

\vspace{0.5em}

\textbf{Aufgabe 3: Verben auf -ar (4 P.)}

el grupo → \lsg{agrupar}, el motivo → \lsg{motivar}

el trabajo → \lsg{trabajar}, el plan → \lsg{planificar/planear}

\vspace{0.5em}

\textbf{Aufgabe 4: Wortfamilien (6 P.)}

educar → \lsg{educación} → educativo

\lsg{participar} → participación → participativo

administrar → \lsg{administración} → \lsg{administrativo}

\vspace{0.5em}

\textbf{Aufgabe 5: Eigenschaften eines guten Chefs (8 P.)}

\lsg{Un buen jefe necesita...}
\begin{enumerate}[leftmargin=2em, nosep]
    \item \lsg{responsabilidad}
    \item \lsg{capacidad de comunicación}
    \item \lsg{flexibilidad}
    \item \lsg{creatividad}
    \item \lsg{capacidad de motivar}
\end{enumerate}

\lsg{Begründung: Un buen jefe necesita flexibilidad porque tiene que adaptarse a diferentes situaciones y personas.}

\end{gruppeBbox}

\end{document}
